% Generated draft: requires mathematical and source review.
\documentclass[11pt]{article}
\usepackage[margin=25mm]{geometry}
\usepackage{fontspec}
\setmainfont{Latin Modern Roman}
\usepackage{amsmath,amssymb,mathtools,unicode-math}
\setmathfont{Latin Modern Math}
\usepackage{enumitem,needspace,xurl,hyperref}
\setlength{\parindent}{0pt}
\setlength{\parskip}{5pt}
\setlength{\emergencystretch}{3em}
\newcommand{\N}{\mathbb{N}}
\newcommand{\Z}{\mathbb{Z}}
\newcommand{\Q}{\mathbb{Q}}
\newcommand{\R}{\mathbb{R}}
\newcommand{\C}{\mathbb{C}}
\newcommand{\F}{\mathbb{F}}
\newcommand{\A}{\mathbb{A}}
\newcommand{\PP}{\mathbb{P}}
\newcommand{\E}{\mathbb{E}}
\newcommand{\eps}{\varepsilon}
\newcommand{\dd}{\,\mathrm{d}}
\begin{document}
\section*{UnsolvedMath ID: 10600047}
\textbf{Source problem code:} catalogue entry
\subsection{Definitions and mathematical statement}
In a usual diagrammatic model, a free knot is represented by an undecorated chord diagram, or equivalently a framed four-valent graph with one unicursal component, modulo the corresponding free Reidemeister moves. The excerpt does not give the precise move convention or specify closed, based, or long free knots. It likewise does not define a connected-sum operation or the connecting choices.

\emph{Three problems.}
\begin{enumerate}
\item Is it true that any connected sum of free knots is trivial? This is the literal supplied question: for each pair and every admissible connected-sum choice, is the result the trivial free knot? The source suggests an affirmative answer and proposes adapting an analogous proof for virtual knots, but the reference is missing. No omitted hypotheses on the factors are supplied.
\item Are free knots algorithmically recognizable? In an equivalence-decision interpretation this asks for an algorithm that halts on every pair of finite free-knot diagrams and decides whether they represent the same free knot. If only recognition of the trivial knot is intended, that is a distinct decision problem. The source suggests a negative answer and speculates that free knots can model many logical constructions; it provides no undecidability theorem or exact recognition convention.
\item Prove that free knots do not commute in general. Once the intended ordered connected sum is defined, a usual noncommutativity target is to find \(K,L\) with
\[
 K*L\not\cong L*K.
\]
The source says standard free-knot methods, perhaps extended, should solve this problem. It does not specify whether the operation is on long, based, or closed diagrams, or how its choices are compared.
\end{enumerate}
The first and third requests cannot both hold for a single well-defined operation on the same equivalence classes if the first says that every product is the trivial class: that would force all products to commute. This is an unresolved conflict in the supplied wording, potentially involving omitted conditions or different categories; it is not silently repaired.

The source contrasts these questions with its claim that flat virtual knots are easily algorithmically recognizable, citing an absent reference, and motivates the difficulty by the then-absent geometric approach to defining free knots. These are supplied contextual assessments, not independently verified claims about present-day algorithms or available geometric methods.
\subsection{Short English statement}
The source asks whether all connected sums are trivial, whether free knots admit an algorithmic recognition procedure, and whether connected sum can be noncommutative. Its first and third requests conflict if read for the same operation and equivalence classes, so the missing category or qualification must be recovered.
\subsection{Sources}
\begin{itemize}
\item Local UnsolvedMath registry, permanent ID 10600047, source code catalogue entry, virtual-knot problem 47. No individual proposer is supplied.
\item \emph{Unsolved Problems in Virtual Knot Theory and Combinatorial Knot Theory}, revised survey, problem-list item 47. Supplied URL: \url{https://arxiv.org/abs/1409.2823}; recorded access 2026-07-29; extraction source-tex. The dedicated section contains 60 top-level enumerated items with compound items retained intact.
\item References for algorithmic recognition of flat virtual knots and the analogous virtual-knot connected-sum proof are missing from the extracted statement.
\item Supplied status partially\_solved and PARTIAL-PROGRESS. The summary says free-knot invariant theory is well developed and these three questions appear open, adding that noncommutativity is likely accessible. The source itself suggests affirmative and negative answers to the first and second questions respectively. These are attributed assessments, not proved results. Current primary-literature status and redistribution terms for the publicly accessible author/publisher copy remain NEEDS\_REVIEW.
\end{itemize}
\end{document}
